% GENERATED FILE. Edit canonical lessons, then run npm run generate:books.
% canonical-source-hash: fnv1a64:8a8e7a3d2b83c7d0
% canonical-lessons: BN-C92-shonibar, BN-C92-robibar, BN-C92-protidin, BN-C92-roj, BN-C92-shit

\chapter{Saturday, Sunday, Every day, Daily, Winter}
\label{ch:bn-saturday-sunday-every-day-daily-winter}
\hlchaptermodality{92}

\begin{chapteropening}
\textbf{By the end of this chapter:} \emph{I can say Saturday, Sunday, every day, daily and winter.}

Complete the last lesson of chapter 92: I can say Saturday, Sunday, every day, daily and winter.
\end{chapteropening}

\section[\texorpdfstring{\bn{শনিবার} (shônibār)}{শনিবার (shônibār)}]{\bn{শনিবার} (shônibār) — Saturday}

\label{lesson:BN-C92-shonibar}



\begin{quote}
Before the new one: say the Bengali for Thursday, then the Bengali for Friday.
\end{quote}

\subsection*{You'll want to know: \bn{শনিবার}}
\textbf{\bn{শনিবার}} — \emph{shônibār} — \textquotedblleft{}Saturday\textquotedblright{}.

\begin{cousinweb}
\textbf{\bn{শনিবার}} (\emph{shônibār}) is Sanskrit \emph{śanivāra}, \textquotedblleft{}Saturday\textquotedblright{}, taken over whole as a learned loan, a tatsama. It is \emph{śani}, \textquotedblleft{}Saturn\textquotedblright{}, + \emph{vāra}, \textquotedblleft{}day of the week\textquotedblright{}.
\end{cousinweb}

\begin{grammarlens}[title={What you've built}]
One more word for this chapter.
\end{grammarlens}

\subsection*{Your turn}
\begin{itemize}
\raggedright
  \item \emph{Say it:} \emph{shônibār}
  \item \emph{Say it:} \emph{shônibār}, once more
  \item \emph{Say it:} say \emph{shônibār}
  \item \emph{Recall:} say the Bengali for Thursday, then the Bengali for Friday, then say \emph{shônibār} again
\end{itemize}

\begin{tcolorbox}[breakable,colback=teal!4,colframe=teal!35!black,title={Before you move on}]
What does \bn{শনিবার} mean? (\textquotedblleft{}Saturday\textquotedblright{}.) Say it once more.
\end{tcolorbox}
\section[\texorpdfstring{\bn{রবিবার} (rôbibār)}{রবিবার (rôbibār)}]{\bn{রবিবার} (rôbibār) — Sunday}

\label{lesson:BN-C92-robibar}



\begin{quote}
Before the new one: say the Bengali for Friday, then the Bengali for Saturday.
\end{quote}

\subsection*{You'll want to know: \bn{রবিবার}}
\textbf{\bn{রবিবার}} — \emph{rôbibār} — \textquotedblleft{}Sunday\textquotedblright{}.

\begin{cousinweb}
\textbf{\bn{রবিবার}} (\emph{rôbibār}) is Sanskrit \emph{ravivāra}, \textquotedblleft{}Sunday\textquotedblright{}, taken over whole as a learned loan, a tatsama. It is \emph{ravi}, \textquotedblleft{}the sun\textquotedblright{}, + \emph{vāra}, \textquotedblleft{}day of the week\textquotedblright{}.
\end{cousinweb}

\begin{grammarlens}[title={What you've built}]
One more word for this chapter.
\end{grammarlens}

\subsection*{Your turn}
\begin{itemize}
\raggedright
  \item \emph{Say it:} \emph{rôbibār}
  \item \emph{Say it:} \emph{rôbibār}, once more
  \item \emph{Say it:} say \emph{rôbibār}
  \item \emph{Recall:} say the Bengali for Friday, then the Bengali for Saturday, then say \emph{rôbibār} again
\end{itemize}

\begin{tcolorbox}[breakable,colback=teal!4,colframe=teal!35!black,title={Before you move on}]
What does \bn{রবিবার} mean? (\textquotedblleft{}Sunday\textquotedblright{}.) Say it once more.
\end{tcolorbox}
\section[\texorpdfstring{\bn{প্রতিদিন} (prôtidin)}{প্রতিদিন (prôtidin)}]{\bn{প্রতিদিন} (prôtidin) — every day}

\label{lesson:BN-C92-protidin}



\begin{quote}
Before the new one: say the Bengali for Saturday, then the Bengali for Sunday.
\end{quote}

\subsection*{You'll want to know: \bn{প্রতিদিন}}
\textbf{\bn{প্রতিদিন}} — \emph{prôtidin} — \textquotedblleft{}every day\textquotedblright{}.

\begin{cousinweb}
\textbf{\bn{প্রতিদিন}} (\emph{prôtidin}) is Sanskrit \emph{pratidinam}, \textquotedblleft{}every day\textquotedblright{}, taken over whole as a learned loan, a tatsama. It is \emph{prati}, \textquotedblleft{}each\textquotedblright{}, + \emph{dina}, \textquotedblleft{}day\textquotedblright{}.
\end{cousinweb}

\begin{grammarlens}[title={What you've built}]
One more word for this chapter.
\end{grammarlens}

\subsection*{Your turn}
\begin{itemize}
\raggedright
  \item \emph{Say it:} \emph{prôtidin}
  \item \emph{Say it:} \emph{prôtidin}, once more
  \item \emph{Say it:} say \emph{prôtidin}
  \item \emph{Recall:} say the Bengali for Saturday, then the Bengali for Sunday, then say \emph{prôtidin} again
\end{itemize}

\begin{tcolorbox}[breakable,colback=teal!4,colframe=teal!35!black,title={Before you move on}]
What does \bn{প্রতিদিন} mean? (\textquotedblleft{}Every day\textquotedblright{}.) Say it once more.
\end{tcolorbox}
\section[\texorpdfstring{\bn{রোজ} (roj)}{রোজ (roj)}]{\bn{রোজ} (roj) — daily}

\label{lesson:BN-C92-roj}



\begin{quote}
Before the new one: say the Bengali for Sunday, then the Bengali for every day.
\end{quote}

\subsection*{You'll want to know: \bn{রোজ}}
\textbf{\bn{রোজ}} — \emph{roj} — \textquotedblleft{}daily\textquotedblright{}.

\begin{cousinweb}
\textbf{\bn{রোজ}} (\emph{roj}) is a loan, usually traced to Persian \emph{rōz}, \textquotedblleft{}day\textquotedblright{}. In Bengali it means \textquotedblleft{}daily\textquotedblright{}.
\end{cousinweb}

\begin{grammarlens}[title={What you've built}]
One more word for this chapter.
\end{grammarlens}

\subsection*{Your turn}
\begin{itemize}
\raggedright
  \item \emph{Say it:} \emph{roj}
  \item \emph{Say it:} \emph{roj}, once more
  \item \emph{Say it:} say \emph{roj}
  \item \emph{Recall:} say the Bengali for Sunday, then the Bengali for every day, then say \emph{roj} again
\end{itemize}

\begin{tcolorbox}[breakable,colback=teal!4,colframe=teal!35!black,title={Before you move on}]
What does \bn{রোজ} mean? (\textquotedblleft{}Daily\textquotedblright{}.) Say it once more.
\end{tcolorbox}
\section[\texorpdfstring{\bn{শীত} (shīt)}{শীত (shīt)}]{\bn{শীত} (shīt) — winter; cold}

\label{lesson:BN-C92-shit}



\begin{quote}
Before the new one: say the Bengali for every day, then the Bengali for daily.
\end{quote}

\subsection*{You'll want to know: \bn{শীত}}
\textbf{\bn{শীত}} — \emph{shīt} — \textquotedblleft{}winter; cold\textquotedblright{}.

\begin{cousinweb}
\textbf{\bn{শীত}} (\emph{shīt}) is Sanskrit \emph{śīta}, \textquotedblleft{}cold\textquotedblright{}, taken over whole as a learned loan, a tatsama.
\end{cousinweb}

\begin{grammarlens}[title={What you've built}]
One more word for this chapter.
\end{grammarlens}

\subsection*{Your turn}
\begin{itemize}
\raggedright
  \item \emph{Say it:} \emph{shīt}
  \item \emph{Say it:} \emph{shīt}, once more
  \item \emph{Say it:} say \emph{shīt}
  \item \emph{Recall:} say the Bengali for every day, then the Bengali for daily, then say \emph{shīt} again
\end{itemize}

\begin{tcolorbox}[breakable,colback=teal!4,colframe=teal!35!black,title={Before you move on}]
What does \bn{শীত} mean? (\textquotedblleft{}Winter; cold\textquotedblright{}.) Say it once more.
\end{tcolorbox}
